\section{Sizes of coefficients and of linear forms}\label{sec:estimates}

We now supply the two size bounds in Proposition~\ref{prop:criterion}.
Both include the multiplier from Section~\ref{sec:arithmetic}: the forms
must remain small after their coefficients have been made integral.
We use derivatives of order exceeding the degree of the finite correction,
while leaving enough decay in the tail for its sum to converge. Accordingly,
assume
\begin{equation}\label{eq:kappa-range}
 2r<\kappa<\omega,\qquad
 2M+2\leq k\leq\kappa n,\qquad 0\leq p\leq h.
\end{equation}
In addition to the integer conditions in Proposition~\ref{prop:rational},
we require $\kappa n$ to be integral. We call such $n$ \emph{admissible}:
$rn$, $\omega n$, $\Omega n$ and $\kappa n$ are all integers.
We prove bounds uniform throughout the displayed range.

\begin{proposition}\label{prop:estimates}
Every coefficient of the integer form $\Lambda_{n,p,k}$ is at most
$B^{n+o(n)}$ in absolute value, where
\begin{equation}\label{eq:B}
 \log B=\log\chi+(r+1)\log(3/2)
                 +\kappa(\log8+\Ht_{a+h}).
\end{equation}
The values satisfy $|\Lambda_{n,p,k}|\leq\alpha^{n+o(n)}$, where
\begin{equation}\label{eq:alpha}
 \log\alpha=\log\chi-\Omega\log r
                 +\kappa(1+\log4+\Ht_{a+h}).
\end{equation}
\end{proposition}

The term $-\Omega\log r$ in \eqref{eq:alpha} supplies the decay;
the other terms are costs from the coefficient bound, differentiation
and denominators. Section~\ref{sec:conclusion} chooses the parameters so
that $\log\alpha<0$ and compares this decay with $\log B$.

\begin{proof}[Coefficient bound]
On the disk $|z+1|\leq1/2$, the row-sum norm of the matrix in
\eqref{eq:base-system} is at most $11/3$, since $|z|\geq1/2$ and
$|1-z|\geq3/2$. Integrating along a straight segment from $-1$ gives
$\|U(z+1)\|\leq e^{11/6}$.
Equations \eqref{eq:coeff-small} and \eqref{eq:C} give
\[
 \max_i |C_i^{[p]}(z)|
     \leq \exp(n\log\chi+o(n))(3/2)^{(r+1)n}
\]
on the same disk. The factors $(i)_p$ depend only on the fixed parameters,
and the number of coefficients is polynomial in $n$.
Cauchy's coefficient estimate therefore bounds
$Q_{k,i}^{[p]}(-1)/(k-1)!$ by this expression times $2^{k-1}$ and a
fixed constant. Multiply by $\Dt_{d,k-1}$.
By monotonicity it is enough to bound $\Dt_{d,\kappa n}$, and
\eqref{eq:D-rate} yields \eqref{eq:B}. The factors $2^{i+1}$ in
\eqref{eq:H-values} are fixed while $n$ tends to infinity.
\end{proof}

\begin{proof}[Value bound]
For real $x\geq M+1/2$ and $0\leq p\leq h$, we first claim
\begin{equation}\label{eq:F-tail-bound}
 |F_n^{(p)}(\pm x)|
 \leq\exp(n\log\chi+o(n))\,r^{-\Omega n}(M/x)^{\omega n}.
\end{equation}
Here $\exp(n\log\chi+o(n))$ is the coefficient cost and
$r^{-\Omega n}$ is the extra smallness supplied by the Laurent bounds.
The remaining factor $(M/x)^{\omega n}$ controls decay as $x$ grows;
it leaves enough power decay to sum after up to $\kappa n-1$
derivatives because $\kappa<\omega$.
The Laurent series converges for $|t|>n$, because all poles lie among
$0,-1,\ldots,-n$. Differentiation adds the factor
$(\nu)_p x^{-p}$ to the bound for the $\nu$th term.
For $\omega n\leq\nu<\Omega n$, combine \eqref{eq:laurent-bounds} as
\[
 r^{\nu-\Omega n}(n/x)^\nu
       =r^{-\Omega n}(M/x)^\nu.
\]
There are $O(n)$ terms, with remaining factors bounded by a fixed power
of $n$. For $\nu\geq\Omega n$, write
$(n/x)^\nu=r^{-\nu}(M/x)^\nu$.
The remaining tail is bounded by a polynomial in $n$ times
$r^{-\Omega n}$, because $r\geq2$ and $a,p$ are fixed.
This proves \eqref{eq:F-tail-bound}, uniformly in $x$ and $p$.

Write $\ell=k-1$.
After division by $\ell!$, the absolute derivatives at $-1$ of the
monomials in \eqref{eq:S} are
\[
 \binom{t-M+\ell}{\ell},
 \qquad
 \binom{M+t}{\ell},
\]
where the second is zero if $M+t<\ell$.
For $x=t+1/2$ and $\ell\geq2M+1$, both upper arguments are at most
$(\ell/M)x$.
The elementary bound $\binom{v}{\ell}\leq(ev/\ell)^\ell$ gives
\[
 \left|\frac{S_{n,p}^{(\ell)}(-1)}{\ell!}\right|
 \leq \exp(n\log\chi+o(n))r^{-\Omega n}e^\ell
       \sum_{t=M}^{\infty}
           \left(\frac{M}{t+1/2}\right)^{\omega n-\ell}.
\]
The exponent $E=\omega n-\ell$ satisfies
$E-1\geq(\omega-\kappa)n$. The sum is at most
$1+M/(E-1)$ by comparison with an integral, so it is uniformly bounded.
We obtain
\[
 \left|\frac{S_{n,p}^{(\ell)}(-1)}{\ell!}\right|
 \leq\exp\bigl(n(\log\chi-\Omega\log r+\kappa)+o(n)\bigr).
\]
These bounds also justify the absolute convergence of the differentiated
series used in Lemma~\ref{lem:tail}. Multiplication by
$\Dt_{d,\ell}\leq\Dt_{d,\kappa n}$ proves \eqref{eq:alpha}.
\end{proof}

